\documentclass[11pt,a4paper]{article}
\usepackage[utf8]{inputenc}
\usepackage[T1]{fontenc}
\usepackage{lmodern}
\usepackage{amsmath,amssymb}
\usepackage[a4paper,top=2.2cm,bottom=2.2cm,left=2.2cm,right=2.2cm,headheight=14pt]{geometry}
\usepackage{enumitem}
\usepackage{tabularx,array}
\usepackage{fancyhdr}
\usepackage{lastpage}
\usepackage{xcolor}
\usepackage{titlesec}
\usepackage{needspace}
\newcolumntype{L}{>{\raggedright\arraybackslash\hyphenpenalty=50}X}
\newcolumntype{P}[1]{>{\raggedright\arraybackslash\hyphenpenalty=50}p{#1}}

\definecolor{akzent}{RGB}{30,80,140}
\setlength{\parindent}{0pt}
\setlength{\parskip}{4pt}
\raggedright
\hyphenpenalty=10000
\exhyphenpenalty=10000
\renewcommand{\arraystretch}{1.25}

\pagestyle{fancy}
\fancyhf{}
\lhead{\small Spezielle Linien im Dreieck}
\rhead{\small Aufgabenblatt}
\cfoot{\small Seite \thepage\ von \pageref{AufgabenEnde}}
\renewcommand{\headrulewidth}{0.4pt}

\titleformat{\section}{\large\bfseries\color{akzent}}{}{0pt}{}
\titlespacing*{\section}{0pt}{14pt}{6pt}
\titleformat{\subsection}{\normalsize\bfseries}{}{0pt}{}
\titlespacing*{\subsection}{0pt}{10pt}{3pt}

\newcommand{\stern}{$\bigstar$\ }
\newcommand{\luecke}{\rule{1.8cm}{0.4pt}}
\setlist[enumerate]{leftmargin=*,itemsep=3pt,topsep=3pt}
\setlist[itemize]{leftmargin=*,itemsep=3pt,topsep=3pt}

\begin{document}

{\LARGE\bfseries Spezielle Linien im Dreieck}\\[2pt]
{\large Repetitions- und Vertiefungsaufgaben}

\vspace{8pt}
Name: \rule{6cm}{0.4pt}\hfill Datum: \rule{3.5cm}{0.4pt}

\vspace{6pt}
\begin{tabularx}{\linewidth}{>{\bfseries}l L r}
\hline
Teil A -- Repetition & Linien, Schnittpunkte und ihre Eigenschaften & ca. 15 Min. \\
Teil B -- Anwenden & Rechnen im Koordinatensystem, Sachaufgaben, L\"angen und Winkel & ca. 25 Min. \\
Teil C -- Vertiefung & Begr\"unden, Lage der Punkte, Euler-Gerade & ca. 20 Min. \\
\hline
\end{tabularx}

\small
Hilfsmittel: Taschenrechner. Es wird nichts konstruiert; l\"ose die Rechnungen auf einem separaten Blatt.
Aufgaben mit \stern sind Zusatzaufgaben f\"ur alle, die schneller fertig sind.\\
Bezeichnungen: $U$ = Umkreismittelpunkt, $I$ = Inkreismittelpunkt, $H$ = H\"ohenschnittpunkt, $S$ = Schwerpunkt.
\normalsize

%==========================================================
\section{Teil A -- Repetition (ca. 15 Min.)}

\subsection{A1 \quad \"Ubersicht vervollst\"andigen}
F\"ulle die Tabelle aus. Die erste Zeile ist als Beispiel bereits ausgef\"ullt.

\vspace{2pt}
\renewcommand{\arraystretch}{1.4}
\begin{tabularx}{\linewidth}{|>{\bfseries}P{3.3cm}|L|P{3.1cm}|L|}
\hline
Linie & \textbf{Definition} & \textbf{Schnittpunkt} & \textbf{Besondere Eigenschaft des Schnittpunkts} \\ \hline
Mittelsenkrechte & steht senkrecht auf einer Seite und geht durch deren Mittelpunkt & Umkreis\-mittelpunkt $U$ & gleich weit von allen drei \textbf{Ecken} entfernt \\ \hline
Winkel\-halbierende & & & \\[1.4cm] \hline
H\"ohe & & & \\[1.4cm] \hline
Seiten\-halbierende & & & \\[1.4cm] \hline
\end{tabularx}
\renewcommand{\arraystretch}{1.25}

\needspace{8cm}
\subsection{A2 \quad Richtig oder falsch?}
Kreuze an und begr\"unde. Korrigiere falsche Aussagen so, dass sie richtig werden.

\begin{tabularx}{\linewidth}{@{}r L c c@{}}
 & & \textbf{R} & \textbf{F} \\
1. & Der Umkreismittelpunkt liegt immer im Innern des Dreiecks. & $\square$ & $\square$ \\[3pt]
2. & Der Inkreismittelpunkt liegt immer im Innern des Dreiecks. & $\square$ & $\square$ \\[3pt]
3. & Der Schwerpunkt teilt jede Seitenhalbierende im Verh\"altnis $1:2$, von der Ecke aus gemessen. & $\square$ & $\square$ \\[3pt]
4. & Jeder Punkt auf der Winkelhalbierenden von $\alpha$ ist von den Seiten $b$ und $c$ gleich weit entfernt. & $\square$ & $\square$ \\[3pt]
5. & In einem rechtwinkligen Dreieck ist der H\"ohenschnittpunkt die Ecke beim rechten Winkel. & $\square$ & $\square$ \\[3pt]
6. & Eine H\"ohe steht senkrecht auf einer Seite und geht durch deren Mittelpunkt. & $\square$ & $\square$ \\[3pt]
7. & Eine Seitenhalbierende teilt das Dreieck in zwei fl\"achengleiche Teildreiecke. & $\square$ & $\square$ \\[3pt]
8. & Der Umkreismittelpunkt ist von allen drei Seiten gleich weit entfernt. & $\square$ & $\square$ \\[3pt]
\end{tabularx}

%==========================================================
\section{Teil B -- Anwenden (ca. 25 Min.)}

\subsection{B1 \quad Rechnen im Koordinatensystem}
Gegeben ist das Dreieck mit $A(0\,|\,1)$, $B(12\,|\,1)$ und $C(3\,|\,10)$; 1 Einheit = 1 cm. L\"ose alles rechnerisch, ohne zu konstruieren.
\begin{enumerate}
\item Berechne die Koordinaten des Schwerpunkts $S$. \emph{Tipp:} Sie sind die Mittelwerte der Eckpunkt-Koordinaten.
\item Begr\"unde, weshalb der Umkreismittelpunkt $U$ auf der Geraden $x = 6$ liegt.
\item Setze $U(6\,|\,y)$. Nutze, dass $U$ von $A$ und $C$ gleich weit entfernt ist (Pythagoras), und berechne $y$ sowie den Umkreisradius.
\item Berechne die H\"ohe $h_c$ und den Fl\"acheninhalt des Dreiecks.
\item Auf welcher Geraden liegt die H\"ohe $h_c$? Was folgt daraus f\"ur die $x$-Koordinate des H\"ohenschnittpunkts $H$?
\item Berechne den Umfang und daraus den Inkreisradius $r$. Es gilt $A = \tfrac12 \cdot r \cdot u$.
\end{enumerate}

\subsection{B2 \quad Sachaufgaben}
Entscheide jeweils, welcher besondere Punkt gesucht ist, und begr\"unde mit seiner Eigenschaft.
\begin{enumerate}
\item \textbf{Feuerwehr:} Drei Gemeinden bauen ein gemeinsames Feuerwehrdepot. Es soll von allen drei Dorfzentren gleich weit entfernt sein. Welches Problem entsteht, wenn die drei Dorfzentren ein stark stumpfwinkliges Dreieck bilden?
\item \textbf{Spielplatz:} Ein dreieckiger Park wird von drei geraden Strassen begrenzt. Darin soll ein m\"oglichst grosser kreisrunder Sandkasten entstehen. Wo liegt sein Mittelpunkt, und wie gross ist er h\"ochstens?
\item \textbf{Balancieren:} Ein d\"unnes dreieckiges Holzbrett soll waagrecht auf einer Bleistiftspitze balancieren. Wo muss die Spitze ansetzen, und wie findest du diesen Punkt nur mit einem Lineal?
\end{enumerate}

\subsection{B3 \quad Berechnungen}
\begin{enumerate}
\item Ein rechtwinkliges Dreieck hat die Katheten 6 cm und 8 cm. Berechne den Umkreisradius. \emph{Tipp:} Satz des Thales.
\item Im selben Dreieck: Berechne den Inkreisradius $r$. \emph{Tipp:} Zerlege das Dreieck in drei Teildreiecke mit gemeinsamer Spitze $I$. Es gilt $A = \tfrac12 \cdot r \cdot u$ ($u$ = Umfang).
\item Die Seitenhalbierende $s_a$ ist 12 cm lang. Wie weit ist $S$ von der Ecke $A$ und vom Seitenmittelpunkt $M_a$ entfernt?
\item In einem Dreieck gilt $\alpha = 70^\circ$ und $\beta = 50^\circ$. Berechne den Winkel $\angle AIB$.
\item \stern Ein Dreieck hat den Fl\"acheninhalt 36 cm$^2$. Die drei Seitenhalbierenden zerlegen es in sechs kleine Dreiecke. Wie gross ist jedes davon? Begr\"unde.
\end{enumerate}

\subsection{B4 \quad L\"angen in besonderen Dreiecken}
\textbf{Gleichschenkliges Dreieck} mit Basis 10 cm und Schenkeln von je 13 cm:
\begin{enumerate}
\item Berechne die H\"ohe auf die Basis und den Fl\"acheninhalt.
\item Wie weit ist der Schwerpunkt $S$ von der Basis entfernt?
\item Berechne den Inkreisradius.
\item \stern Berechne den Umkreisradius $R$. \emph{Tipp:} $U$ liegt auf der Symmetrieachse. Stelle mit Pythagoras eine Gleichung f\"ur $R$ auf.
\end{enumerate}
\textbf{Rechtwinkliges Dreieck} mit den Katheten 5 cm und 12 cm (rechter Winkel bei $C$):
\begin{enumerate}
\item Berechne den Umkreisradius.
\item Wie lang ist die Seitenhalbierende $s_c$ auf die Hypotenuse? \emph{Tipp:} Thaleskreis.
\item Wie weit ist der Schwerpunkt $S$ von der Ecke $C$ entfernt?
\item Berechne den Inkreisradius.
\end{enumerate}

\subsection{B5 \quad Winkel}
Im Dreieck aus B3.4 gilt $\alpha = 70^\circ$ und $\beta = 50^\circ$.
\begin{enumerate}
\item Berechne den Winkel zwischen der H\"ohe $h_a$ und der Seite $c$.
\item Berechne den Winkel $\angle BHC$. \emph{Tipp:} Betrachte das Viereck aus $A$, $H$ und den H\"ohenfusspunkten auf $b$ und $c$.
\item \stern Zeige allgemein: In jedem Dreieck gilt $\angle AIB = 90^\circ + \dfrac{\gamma}{2}$.
\end{enumerate}

%==========================================================
\section{Teil C -- Vertiefung (ca. 20 Min.)}

\subsection{C1 \quad Begr\"unden}
Weshalb schneiden sich die drei Mittelsenkrechten immer in einem einzigen Punkt? Erg\"anze die Begr\"undung:

\begin{itemize}[itemsep=10pt]
\item Jeder Punkt auf der Mittelsenkrechten von $AB$ ist gleich weit von \luecke\ und \luecke\ entfernt.
\item Jeder Punkt auf der Mittelsenkrechten von $BC$ ist gleich weit von \luecke\ und \luecke\ entfernt.
\item Der Schnittpunkt $U$ dieser beiden Linien ist also gleich weit von \rule{4.5cm}{0.4pt} entfernt.
\item Deshalb liegt $U$ auch auf \rule{4.5cm}{0.4pt}, denn \rule{6cm}{0.4pt}\\[10pt] \rule{0.96\linewidth}{0.4pt}.
\end{itemize}

\stern Formuliere eine entsprechende Begr\"undung f\"ur die Winkelhalbierenden.

\needspace{9cm}
\subsection{C2 \quad Lage der Punkte je nach Dreiecksform}
\"Uberlege dir (allenfalls mit einer Freihandskizze), wo die Punkte im spitz-, recht- und stumpfwinkligen Dreieck liegen. Trage ein, ob der Punkt innen, auf dem Rand (wo genau?) oder aussen liegt.

\vspace{2pt}
\renewcommand{\arraystretch}{2.1}
\begin{tabularx}{\linewidth}{|>{\bfseries}c|L|L|L|}
\hline
Punkt & \textbf{spitzwinklig} & \textbf{rechtwinklig} & \textbf{stumpfwinklig} \\ \hline
$U$ & & & \\ \hline
$H$ & & & \\ \hline
$S$ & & & \\ \hline
$I$ & & & \\ \hline
\end{tabularx}
\renewcommand{\arraystretch}{1.25}

\subsection{C3 \quad Besondere Dreiecke}
\begin{enumerate}
\item In welchem Dreieck fallen $U$, $H$, $S$ und $I$ zusammen? Begr\"unde.
\item In einem gleichschenkligen Dreieck liegen alle vier Punkte auf einer gemeinsamen Geraden. Auf welcher, und weshalb?
\end{enumerate}

\subsection{C4 \quad Dreieck r\"uckw\"arts}
Von einem Dreieck sind nur $A(1\,|\,2)$, $B(7\,|\,0)$ und der Schwerpunkt $S(4\,|\,3)$ bekannt. Bestimme die Koordinaten von $C$. \emph{Tipp:} Die Koordinaten von $S$ sind die Mittelwerte der Eckpunkt-Koordinaten.

\subsection{C5 \quad \stern Euler-Gerade}
In jedem Dreieck liegen $U$, $S$ und $H$ auf einer Geraden, der sogenannten Euler-Geraden. Dabei liegt $S$ zwischen $H$ und $U$, und es gilt $\overline{HS} : \overline{SU} = 2:1$.
\begin{enumerate}
\item Berechne mit $S$ und $U$ aus B1 die Koordinaten von $H$. Passt das Resultat zu B1.5?
\item Liegt auch $I$ auf der Euler-Geraden? \emph{Tipp:} $I$ ist von der Seite $AB$ (auf der Geraden $y = 1$) genau $r$ entfernt.
\item Weshalb ist die Euler-Gerade in einem gleichseitigen Dreieck nicht definiert?
\end{enumerate}

\label{AufgabenEnde}
%==========================================================
\newpage
\setcounter{page}{1}
\rhead{\small L\"osungen}
\cfoot{\small L\"osungen, Seite \thepage}
{\LARGE\bfseries L\"osungen}

\section{Teil A}

\subsection{A1}
\small
\begin{tabularx}{\linewidth}{|>{\bfseries}P{3.1cm}|L|P{3cm}|L|}
\hline
Linie & \textbf{Definition} & \textbf{Schnittpunkt} & \textbf{Eigenschaft} \\ \hline
Mittelsenkrechte & senkrecht auf einer Seite, durch deren Mittelpunkt & Umkreis\-mittelpunkt $U$ & gleich weit von allen drei Ecken entfernt (Mittelpunkt des Umkreises) \\ \hline
Winkel\-halbierende & halbiert einen Innenwinkel & Inkreis\-mittelpunkt $I$ & gleich weit von allen drei Seiten entfernt (Mittelpunkt des Inkreises) \\ \hline
H\"ohe & durch eine Ecke, senkrecht auf die gegen\"uberliegende Seite (bzw.\ deren Verl\"angerung) & H\"ohen\-schnittpunkt $H$ & liegt mit $U$ und $S$ auf der Euler-Geraden \\ \hline
Seiten\-halbierende & verbindet eine Ecke mit dem Mittelpunkt der gegen\"uberliegenden Seite & Schwerpunkt $S$ & teilt jede Seitenhalbierende im Verh\"altnis $2:1$ (von der Ecke aus); Gleichgewichtspunkt \\ \hline
\end{tabularx}
\normalsize

\subsection{A2}
\begin{enumerate}
\item Falsch. Beim stumpfwinkligen Dreieck liegt $U$ aussen, beim rechtwinkligen auf der Hypotenusenmitte.
\item Richtig. Alle Winkelhalbierenden verlaufen im Innern des Dreiecks.
\item Falsch. Das Verh\"altnis ist $2:1$, von der Ecke aus gemessen.
\item Richtig. Das ist die Ortslinien-Eigenschaft der Winkelhalbierenden.
\item Richtig. Die beiden Katheten sind zugleich H\"ohen; sie schneiden sich in der Ecke beim rechten Winkel.
\item Falsch. Das beschreibt die Mittelsenkrechte. Eine H\"ohe geht durch die gegen\"uberliegende Ecke.
\item Richtig. Beide Teildreiecke haben gleich lange Grundseiten (halbe Seite) und dieselbe H\"ohe.
\item Falsch. $U$ ist von den drei \textbf{Ecken} gleich weit entfernt, $I$ von den drei Seiten.
\end{enumerate}

\section{Teil B}

\subsection{B1}
\begin{enumerate}
\item $S = \left(\frac{0+12+3}{3}\,\middle|\,\frac{1+1+10}{3}\right) = \mathbf{(5\,|\,4)}$.
\item $A$ und $B$ haben dieselbe $y$-Koordinate, $AB$ ist also waagrecht. Die Mittelsenkrechte von $AB$ steht senkrecht dazu und geht durch den Mittelpunkt $(6\,|\,1)$, ist also die Gerade $x = 6$.
\item $6^2 + (y-1)^2 = 3^2 + (y-10)^2 \;\Rightarrow\; 36 + y^2 - 2y + 1 = 9 + y^2 - 20y + 100 \;\Rightarrow\; 18y = 72 \;\Rightarrow\; y = 4$.\\
$U = \mathbf{(6\,|\,4)}$, Umkreisradius $= \sqrt{36+9} = \sqrt{45} \approx \mathbf{6.71}$ \textbf{cm}.
\item Weil $AB$ waagrecht ist: $h_c = 10 - 1 = \mathbf{9}$ \textbf{cm}, $A = \tfrac12 \cdot 12 \cdot 9 = \mathbf{54}$ \textbf{cm}$^2$.
\item $h_c$ steht senkrecht auf $AB$ und geht durch $C$, liegt also auf der Geraden $x = 3$. Damit hat $H$ die $x$-Koordinate 3.
\item $a = \sqrt{9^2+9^2} = \sqrt{162} \approx 12.73$ cm, $b = \sqrt{3^2+9^2} = \sqrt{90} \approx 9.49$ cm, $c = 12$ cm.\\
$u \approx 34.21$ cm, \ $r = \dfrac{2A}{u} = \dfrac{108}{34.21} \approx \mathbf{3.16}$ \textbf{cm}.
\end{enumerate}

\subsection{B2}
\begin{enumerate}
\item Umkreismittelpunkt $U$: gleich weit von allen drei Ecken. Bei einem stark stumpfwinkligen Dreieck liegt $U$ ausserhalb und weit weg; das Depot st\"unde ausserhalb des Gebiets, alle Wege w\"aren unn\"otig lang.
\item Inkreismittelpunkt $I$: gleich weit von allen drei Seiten (Strassen). Der gr\"osstm\"ogliche Sandkasten ist der Inkreis; sein Radius ist der Inkreisradius.
\item Schwerpunkt $S$. Zwei Seitenmitten mit dem Lineal ausmessen, jeweils mit der gegen\"uberliegenden Ecke verbinden; der Schnittpunkt ist $S$.
\end{enumerate}

\subsection{B3}
\begin{enumerate}
\item Hypotenuse $= 10$ cm. Nach Thales ist sie ein Durchmesser des Umkreises, also Umkreisradius $= \mathbf{5}$ \textbf{cm}.
\item $u = 6+8+10 = 24$ cm, $A = \tfrac12 \cdot 6 \cdot 8 = 24$ cm$^2$. Aus $24 = \tfrac12 \cdot r \cdot 24$ folgt $r = \mathbf{2}$ \textbf{cm}.
\item $\overline{AS} = \tfrac23 \cdot 12 = \mathbf{8}$ \textbf{cm}, \ $\overline{SM_a} = \tfrac13 \cdot 12 = \mathbf{4}$ \textbf{cm}.
\item Im Dreieck $ABI$ sind die Winkel bei $A$ und $B$ halbiert: $35^\circ$ und $25^\circ$. Also $\angle AIB = 180^\circ - 35^\circ - 25^\circ = \mathbf{120^\circ}$.
\item Je \textbf{6 cm}$^\mathbf{2}$. Die zwei Teildreiecke an derselben Seitenmitte haben gleiche Grundseite und gleiche H\"ohe, sind also gleich gross. Da jede Seitenhalbierende das Dreieck halbiert, sind alle drei Paare gleich gross: $36 : 6 = 6$ cm$^2$.
\end{enumerate}

\subsection{B4}
\textbf{Gleichschenkliges Dreieck}
\begin{enumerate}
\item $h = \sqrt{13^2 - 5^2} = \mathbf{12}$ \textbf{cm}, \ $A = \tfrac12 \cdot 10 \cdot 12 = \mathbf{60}$ \textbf{cm}$^\mathbf{2}$.
\item Die H\"ohe auf die Basis ist zugleich Seitenhalbierende. $S$ teilt sie im Verh\"altnis $2:1$, ist also $\tfrac13 \cdot 12 = \mathbf{4}$ \textbf{cm} von der Basis entfernt.
\item $u = 36$ cm, \ $r = \dfrac{2A}{u} = \dfrac{120}{36} = \dfrac{10}{3} \approx \mathbf{3.33}$ \textbf{cm}.
\item Abstand von $U$ zur Basis: $12 - R$. Mit Pythagoras: $R^2 = 5^2 + (12-R)^2 \;\Rightarrow\; 24R = 169 \;\Rightarrow\; R = \dfrac{169}{24} \approx \mathbf{7.04}$ \textbf{cm}.
\end{enumerate}
\textbf{Rechtwinkliges Dreieck}
\begin{enumerate}
\item Hypotenuse $= \sqrt{5^2+12^2} = 13$ cm. Nach Thales ist sie ein Durchmesser, also Umkreisradius $= \mathbf{6.5}$ \textbf{cm}.
\item $s_c$ verbindet $C$ mit dem Mittelpunkt der Hypotenuse, also mit $U$. Darum gilt $s_c$ = Umkreisradius $= \mathbf{6.5}$ \textbf{cm}.
\item $\overline{CS} = \tfrac23 \cdot 6.5 = \tfrac{13}{3} \approx \mathbf{4.33}$ \textbf{cm}.
\item $u = 30$ cm, $A = \tfrac12 \cdot 5 \cdot 12 = 30$ cm$^2$, \ $r = \dfrac{2A}{u} = \mathbf{2}$ \textbf{cm}.
\end{enumerate}

\subsection{B5}
\begin{enumerate}
\item $h_a$, die Seite $c$ und die Strecke von $B$ zum H\"ohenfusspunkt bilden ein rechtwinkliges Dreieck mit dem Winkel $\beta = 50^\circ$ bei $B$. Gesuchter Winkel: $90^\circ - 50^\circ = \mathbf{40^\circ}$.
\item Im Viereck aus $A$, den beiden H\"ohenfusspunkten und $H$ sind zwei Winkel $90^\circ$. Winkel bei $H$: $360^\circ - 90^\circ - 90^\circ - 70^\circ = 110^\circ$. $\angle BHC$ ist sein Scheitelwinkel: $\mathbf{110^\circ}$ (allgemein $180^\circ - \alpha$).
\item $\angle AIB = 180^\circ - \dfrac{\alpha}{2} - \dfrac{\beta}{2} = 180^\circ - \dfrac{\alpha+\beta}{2} = 180^\circ - \dfrac{180^\circ - \gamma}{2} = \mathbf{90^\circ + \dfrac{\gamma}{2}}$. Kontrolle mit B3.4: $90^\circ + 30^\circ = 120^\circ$.
\end{enumerate}

\section{Teil C}

\subsection{C1}
\ldots\ von $A$ und $B$ \ldots\ von $B$ und $C$ \ldots\ von $A$, $B$ und $C$ \ldots\ auf der Mittelsenkrechten von $AC$, denn er ist gleich weit von $A$ und $C$ entfernt.

\stern Ein Punkt auf $w_\alpha$ ist gleich weit von $b$ und $c$ entfernt, einer auf $w_\beta$ gleich weit von $a$ und $c$. Der Schnittpunkt ist also gleich weit von $a$ und $b$ entfernt und liegt deshalb auch auf $w_\gamma$.

\subsection{C2}
\begin{tabularx}{\linewidth}{|>{\bfseries}c|L|L|L|}
\hline
Punkt & \textbf{spitzwinklig} & \textbf{rechtwinklig} & \textbf{stumpfwinklig} \\ \hline
$U$ & innen & Mittelpunkt der Hypotenuse & aussen \\ \hline
$H$ & innen & Ecke beim rechten Winkel & aussen \\ \hline
$S$ & innen & innen & innen \\ \hline
$I$ & innen & innen & innen \\ \hline
\end{tabularx}

\subsection{C3}
\begin{enumerate}
\item Im gleichseitigen Dreieck. Dort fallen Mittelsenkrechte, Winkelhalbierende, H\"ohe und Seitenhalbierende zu jeder Seite zusammen (Symmetrieachsen), also auch ihre Schnittpunkte.
\item Auf der Symmetrieachse. Zur Basis fallen Mittelsenkrechte, H\"ohe, Seitenhalbierende und die Winkelhalbierende des Winkels an der Spitze zusammen.
\end{enumerate}

\subsection{C4}
$x = 3 \cdot 4 - 1 - 7 = 4$, \ $y = 3 \cdot 3 - 2 - 0 = 7$, also $\mathbf{C(4\,|\,7)}$.

\subsection{C5}
\begin{enumerate}
\item Von $U$ nach $S$ geht man um $(-1\,|\,0)$. Weil $\overline{HS}$ doppelt so lang ist wie $\overline{SU}$, gilt $H = S + 2 \cdot (-1\,|\,0) = \mathbf{(3\,|\,4)}$. Das passt zu B1.5: $x = 3$.
\item Nein. $I$ hat die $y$-Koordinate $1 + r \approx 4.16$, liegt also nicht auf der Euler-Geraden $y = 4$. Nur in gleichschenkligen Dreiecken liegt $I$ im Allgemeinen auf der Euler-Geraden.
\item $U$, $S$ und $H$ fallen in einem Punkt zusammen. Durch einen einzigen Punkt ist keine Gerade eindeutig festgelegt.
\end{enumerate}

\end{document}
